\documentclass[11pt]{article}
\usepackage[T1]{fontenc}
\usepackage{lmodern}
\usepackage[margin=1in]{geometry}
\usepackage{amsmath,amssymb,amsthm,mathtools,booktabs,microtype,xcolor}
\usepackage[colorlinks=true,linkcolor=blue!45!black,urlcolor=blue!55!black]{hyperref}
\newtheorem{theorem}{Theorem}[section]
\newtheorem{proposition}[theorem]{Proposition}
\newtheorem{lemma}[theorem]{Lemma}
\newtheorem{corollary}[theorem]{Corollary}
\theoremstyle{definition}
\newtheorem{definition}[theorem]{Definition}
\newtheorem{example}[theorem]{Example}
\theoremstyle{remark}
\newtheorem{remark}[theorem]{Remark}
\newcommand{\F}{\mathbb F}
\newcommand{\R}{\mathbb R}
\newcommand{\C}{\mathbb C}
\newcommand{\PP}{\mathbb P}
\newcommand{\PGL}{\operatorname{PGL}}
\newcommand{\im}{\operatorname{im}}
\newcommand{\id}{\operatorname{id}}
\newcommand{\ii}{\mathrm i}
\title{Conic Motions, Arithmetic Histories, and Three Finite Calculi}
\author{Research note in the math-notes programme}
\date{November 6, 2026}
\begin{document}
\maketitle
\begin{abstract}
The eight-point circle over $\F_7$ is placed inside the projective arithmetic
semantics of AEG. Its one-step rotation becomes a four-context arithmetic
history with an eight-cycle on $\PP^1(\F_7)$. We prove the distinction between
its projective relation and its ordinary evaluation domains, and explain why
directly projectivizing norm-one multiplication loses a factor of two.
Split, nonsplit, and unipotent motions give six-, eight-, and seven-cycles.
Their differences have the same elementary solvability condition, but sharply
different spectral and normalization properties in characteristic seven.
All statements have proofs; an independent exact verification script checks
the finite examples. The proposed Paper I interface is separated from the
open problem of an elliptic arithmetic expression surface.
\end{abstract}
\tableofcontents

\section{Scope, choices, and conventions}
This note continues \href{48-finite-field-trigonometry.pdf}{companion note 48 on finite trigonometry}. Its immediate
target is an explicit example in Paper I's existing projective semantics,
not a change to the definition of a regular arithmetic expression space
(AES). An elliptic projective motion, an affine conic of elliptic type, and
a positive-curvature AES are different objects.

All substantial results below are \emph{proved with stated hypotheses} in
this note. The finite computations are supplementary checks, not proofs for
arbitrary fields. Statements described as an interface or a next problem are
proposals, not changes to the AEG manuscripts or their status registers.

Matrices act on columns from the left, with affine coordinate $t=U/V$:
\[
 \begin{pmatrix}a&b\\c&d\end{pmatrix}[U:V]
   =[aU+bV:cU+dV].
\]
For a chronological history $(g_1,\ldots,g_m)$, the operator is
$g_m\circ\cdots\circ g_1$. Write
$A_b(t)=t+b$, $M_a(t)=at$ for $a\ne0$, and $J(t)=1/t$.
Here $J$ is positive reciprocal, not the negative inversion used in some
Hecke conventions. Projectively these are everywhere defined on $\PP^1$;
ordinary evaluation of $J$ requires a nonzero finite input.

For an affine conic one must retain a chosen line $L_\infty$.
For the present circle, the data $(C,L_\infty,P_0,\tau)$ retain the affine
structure, a base point, and a selected one-step rotation. Two nonrational
geometric intersection points of $C$ with $L_\infty$ give the nonsplit case;
two rational points give the split case; tangency gives the parabolic case.
The bare projective conic does not select a difference operator or a cyclic
ordering of its rational points.

\section{From circle multiplication to a projective torus}
Let $K$ be a field of characteristic different from two in which $-1$ is not
a square, and let $L=K[\ii]/(\ii^2+1)$, with conjugation $\bar\ii=-\ii$.
The conic $C:X^2+Y^2=Z^2$ has parametrization
\begin{equation}\label{eq:nu}
 \nu([U:V])=[V^2-U^2:2UV:U^2+V^2].
\end{equation}
There are no rational points on $C\cap\{Z=0\}$. All rational points of $C$
are therefore affine, and in the finite chart
\[
 w(t)=x(t)+\ii y(t)=\frac{1+\ii t}{1-\ii t},\qquad
 t=\frac{y}{1+x}.
\]
The exceptional point is $\nu(\infty)=(-1,0)$. The homogeneous formula
\eqref{eq:nu} is an isomorphism of the projective conic with $\PP^1$:
away from $(-1,0)$ its inverse is $[U:V]=[Y:Z+X]$;
near that point use $[U:V]=[Z-X:Y]$. These agree on their overlap by
$Y^2=(Z-X)(Z+X)$.

\begin{proposition}[The quotient and norm-one actions]\label{prop:torus}
There is a group isomorphism
\[
 \rho:L^\times/K^\times\longrightarrow L^1,
 \qquad [\beta]\longmapsto\beta/\bar\beta,
 \qquad L^1=\{w:w\bar w=1\}.
\]
For $\beta=a+b\ii$, the corresponding action on $\PP^1(K)$ is
\begin{equation}\label{eq:beta-action}
 t\longmapsto\frac{at+b}{a-bt},\qquad
 B_\beta=\begin{pmatrix}a&b\\-b&a\end{pmatrix}.
\end{equation}
The parametrization intertwines this action with multiplication by
$\rho([\beta])$ on $C(K)=L^1$.
\end{proposition}
\begin{proof}
The determinant of $B_\beta$ is $a^2+b^2\ne0$. Multiplication of these
matrices is multiplication of the corresponding $\beta$'s; scalar matrices
are exactly $b=0$. The kernel of $\beta\mapsto\beta/\bar\beta$ is
$K^\times$. For any $w\in L^1$ with $w\ne-1$, choose $\beta=1+w$:
$\bar\beta=(1+w)/w$, so $\beta/\bar\beta=w$.
For $w=-1$, choose $\beta=\ii$. Thus the map is surjective as well.
Finally direct substitution gives
\[
 \frac{1+\ii(at+b)/(a-bt)}{1-\ii(at+b)/(a-bt)}
   =\frac{a+b\ii}{a-b\ii}\frac{1+\ii t}{1-\ii t}.
\]
The equality is homogeneous and hence also holds at exceptional chart points.
\end{proof}

\begin{remark}[A factor of two that must be retained]\label{rem:double}
Direct multiplication by a norm-one element $w$ on the two-dimensional
$K$-vector space $L$, followed by projectivization, acts on direction
classes $[z]\mapsto[wz]$. Under $\rho$ this becomes
$\rho([wz])=w^2\rho([z])$.
Thus it corresponds to the \emph{square} of the circle rotation by $w$.
Over $\F_7$, $\alpha=2+2\ii$ has order eight, but its direct multiplication
matrix has projective order four because $\alpha^4=-1$ is scalar.
This does not contradict Proposition~\ref{prop:torus}: to obtain rotation by
$\alpha$, use a class $[\beta]$ with $\beta/\bar\beta=\alpha$.
\end{remark}

\section{An elliptic history and its ordinary domains}
From now through Section~\ref{sec:covariance}, take $K=\F_7$ unless another
coefficient field is explicitly specified. For $\beta=1+3\ii$,
\[
 \beta/\bar\beta=2+2\ii=\alpha,\qquad
 R(t)=\frac{t+3}{1-3t},\qquad
 R=\left[\begin{pmatrix}1&3\\4&1\end{pmatrix}\right].
\]
The determinant is $3\ne0$. The stereographic coordinates of $\alpha^k$ are
\begin{equation}\label{eq:cycle}
 (t_0,\ldots,t_7)=(0,3,1,5,\infty,2,6,4),\qquad R(t_k)=t_{k+1}.
\end{equation}
Indices in this formula are modulo eight.

\begin{theorem}[A finite elliptic arithmetic-history witness]
\label{thm:history-witness}
Let $W=(A_2,J,M_2,A_2)$ be a chronological history over $\F_7$.
Its projective operator is $R$, of order eight. For $m\ge0$, let $D_m$
be the set of finite initial values at which the repeated word $W^m$ is
ordinarily admissible. Then
\[
 D_m=\{t_j:j\in\{0,1,2,3,5,6,7\},\quad
             j+r\not\equiv3\pmod8\ (0\le r<m)\},
 \qquad |D_m|=\max(7-m,0).
\]
In particular $W^8$ and the empty word have the same projective operator,
while their ordinary domains are respectively $\varnothing$ and $\F_7$.
\end{theorem}
\begin{proof}
On its ordinary domain,
\[
 W(t)=2+\frac2{t+2}=\frac{2t+6}{t+2}
      =\frac{t+3}{1-3t}.
\]
Its matrix is $\left(\begin{smallmatrix}2&6\\1&2\end{smallmatrix}\right)$,
twice the displayed matrix for $R$. The only ordinary exclusion for one
copy is $t=5=t_3$, where the reciprocal divides by zero. Equation
\eqref{eq:cycle} follows by evaluating $R$ on the eight projective points;
as it is one eight-cycle, $R$ has order eight. Each iteration fails exactly
when its initial state is $t_3$. For the seven finite initial states, the
first occurrence of this failure has index $r=0,1,\ldots,6$, each once;
the remaining index $r=7$ belongs to the inadmissible initial state
$t_4=\infty$. This proves the domain formula and count. The last assertion
compares these domains with the projective identity $R^8=1$.
\end{proof}

Thus ordinary domains become empty after seven copies already. Eight copies
are distinguished because that is the first positive projective identity
relation. The witness concerns a particular marked arithmetic word, not just
the simplified rational function representing its final operator.

\section{Three motions on the same projective line}\label{sec:three}
Consider $R$ above, $S(t)=3t$, and $U(t)=t+1$. The element $3$ generates
$\F_7^\times$. Their orbit decompositions are
\[
\begin{array}{ccl}
 R&:&(0,3,1,5,\infty,2,6,4),\\
 S&:&(1,3,2,6,4,5),\quad(0),\quad(\infty),\\
 U&:&(0,1,2,3,4,5,6),\quad(\infty).
\end{array}
\]
For a representative $A$, the zero/nonzero square class of
$(\operatorname{tr}A)^2-4\det A$ is unchanged by scalar multiplication.
It is $6$ for $R$, $4$ for $S$, and $0$ for $U$.
Thus these are respectively nonsplit semisimple, split semisimple, and
nontrivial unipotent examples. These terms refer to projective operators;
they do not assert a curvature sign for a two-dimensional AES.

The subgroup in Proposition~\ref{prop:torus} has order
$(49-1)/(7-1)=8$, and $R$ generates it. The scaling subgroup generated by
$S$ has order six. The translation subgroup generated by $U$ has order
seven. These give the three moving orbits on which we compare differences.
For reference, $|\PGL_2(\F_7)|=(49-1)(49-7)/(7-1)=336$:
choose a nonzero first column, a second column outside its span, and then
divide by the six nonzero scalar representatives of each projective matrix.
At each extra fixed point, the difference is zero, so the full point-function
space has one constant kernel vector per orbit.

\section{Cyclic difference, solvability, and products}
Let $E$ be any field, let $n\ge2$, and put
$V_n(E)=\{f:\mathbb Z/n\mathbb Z\to E\}$.
Define
\[
 (Tf)(k)=f(k+1),\quad \Delta=T-I,\quad
 \sigma(f)=\sum_{k=0}^{n-1}f(k).
\]
\begin{proposition}[Periodic antiderivatives]\label{prop:antiderivative}
For every $E$ and $n$,
\[
 \ker\Delta=E\mathbf1,\qquad
 \im\Delta=\ker\sigma.
\]
For $g\in\ker\sigma$, all solutions of $\Delta f=g$ are
\begin{equation}\label{eq:primitive}
 f(k)=c+\sum_{j=0}^{k-1}g(j)\qquad(0\le k<n),\quad c\in E.
\end{equation}
If $n$ is invertible in $E$, there is a unique solution with $\sigma(f)=0$.
If $\operatorname{char}E$ divides $n$, adding a constant does not change
$\sigma(f)$; zero sum cannot select a unique solution and may be impossible.
\end{proposition}
\begin{proof}
The kernel consists of functions constant along the cycle. Summing a
difference telescopes, so the image lies in $\ker\sigma$. Conversely
\eqref{eq:primitive} satisfies the equation, including the closing edge
precisely when $\sigma(g)=0$. Two solutions differ by a constant. Their sums
differ by $nc$, giving the normalization claims.
\end{proof}

\begin{proposition}[Averaging and its obstruction]\label{prop:average}
If $n$ is invertible in $E$, then
\[
 V_n(E)=E\mathbf1\oplus\ker\sigma,
\]
and $\Delta$ restricts to an automorphism of $\ker\sigma$.
If $\operatorname{char}E$ divides $n$, this direct sum fails and there is
no linear right inverse to $\Delta:V_n(E)\to\ker\sigma$ that
commutes with $T$.
\end{proposition}
\begin{proof}
For the first assertion subtract $n^{-1}\sigma(f)\mathbf1$ from $f$.
The restriction has zero kernel and maps a finite-dimensional space to
itself, hence is invertible. In the second case $\mathbf1\in\ker\sigma$.
If a commuting right inverse $B$ existed, then
$T(B\mathbf1)=B(T\mathbf1)=B\mathbf1$; thus $B\mathbf1$ would be constant
and $\Delta B\mathbf1=0$, contradicting the right inverse property.
\end{proof}

\begin{proposition}[Product and summation-by-parts formulas]
\label{prop:products}
For pointwise products,
\[
 \Delta(fg)=(\Delta f)g+(Tf)\Delta g
          =f\Delta g+g\Delta f+(\Delta f)(\Delta g),
\]
and
\[
 \sum_k (\Delta f)(k)g(k)
    =\sum_k f(k)\bigl(g(k-1)-g(k)\bigr).
\]
Every $E$-linear derivation of the point-function algebra $V_n(E)$ is zero.
\end{proposition}
\begin{proof}
Expand the products and reindex the finite sums. For the last assertion,
the point indicators $e_k$ are a basis of orthogonal idempotents. A derivation
$D$ obeys $D(e_k)=2e_kD(e_k)$. At a point outside the support of $e_k$ this
forces $D(e_k)=0$; at its support it gives $D(e_k)=2D(e_k)$, again forcing
zero in every characteristic. Thus $D$ vanishes on the entire basis.
\end{proof}
Consequently an ordinary formal derivative on polynomials or rational
functions is a different object from $\Delta$ on point functions. For example,
$t^7-t$ vanishes on $\F_7$, but its formal derivative is $-1$.

\section{The two semisimple calculi}\label{sec:semisimple}
\begin{theorem}[Fourier diagonalization]\label{thm:fourier}
Suppose $n$ is invertible in $E$ and $E$ contains an element $\zeta$ of order
$n$. Then $\chi_j(k)=\zeta^{jk}$, $0\le j<n$, form a basis of $V_n(E)$,
\[
 \Delta\chi_j=(\zeta^j-1)\chi_j,\qquad
 \widehat f(j)=\sum_k f(k)\zeta^{-jk},\qquad
 f(k)=n^{-1}\sum_j\widehat f(j)\zeta^{jk}.
\]
On the zero-sum subspace, the unique zero-sum antiderivative has Fourier
coefficients $\widehat g(j)/(\zeta^j-1)$ for $j\ne0$, and zero for $j=0$.
\end{theorem}
\begin{proof}
For $j\ne l$, geometric summation gives
$\sum_k\zeta^{(j-l)k}=0$, whereas it is $n$ for $j=l$.
These equalities prove inversion and linear independence. Applying $T$ to
each character proves the eigenvalue formula; only $j=0$ has eigenvalue zero.
Dividing the remaining Fourier coefficients proves the final assertion.
\end{proof}

For the split six-cycle take $E=\F_7$, $\zeta=3$. For the nonsplit eight-cycle
take $E=\F_{49}$, $\zeta=\alpha$. In both cases the rank of every positive
power of $\Delta$ is $n-1$. This rank assertion also holds for $\F_7$-valued
functions on the eight-cycle, since extending scalars preserves matrix rank.
Diagonalization of that eight-cycle requires the larger field, however.

\begin{proposition}[Descent of the elliptic Fourier coefficients]
\label{prop:frobenius}
For the eight-cycle, a function $f:\mathbb Z/8\to\F_{49}$ is $\F_7$-valued
if and only if
\[
 \widehat f(-j)=\widehat f(j)^7\qquad(j\in\mathbb Z/8).
\]
For such a $g$ with zero sum, the Fourier antiderivative in
Theorem~\ref{thm:fourier} is also $\F_7$-valued.
\end{proposition}
\begin{proof}
Since $\alpha^7=\alpha^{-1}$, taking seventh powers of the Fourier formula
gives the stated relation when $f^7=f$. Conversely that relation and Fourier
inversion give $f(k)^7=f(k)$, so every value belongs to $\F_7$.
Seventh powers send $\alpha^j-1$ to $\alpha^{-j}-1$, so dividing nonzero
modes preserves the relation; the zero coefficient is fixed as well.
\end{proof}
Thus a zero-sum $\F_7$-valued function has a unique zero-sum primitive over
$\F_7$ in both semisimple examples, although their diagonalizing fields differ.

\section{The unipotent calculus in characteristic seven}
Identify the seven-cycle with $k\in\F_7$.
\begin{theorem}[A divided-power basis for the seven-cycle]
\label{thm:unipotent}
On $V_7(\F_7)$,
\[
 \Delta^7=0,\qquad \operatorname{rank}\Delta^r=7-r
       \quad(0\le r\le7).
\]
More explicitly, the polynomial point functions
\[
 b_r(k)=\binom{k}{r}
       =\frac{k(k-1)\cdots(k-r+1)}{r!},\quad0\le r\le6,
\]
with $b_0=1$, form a basis and satisfy
$\Delta b_0=0$, $\Delta b_r=b_{r-1}$ for $1\le r\le6$.
Consequently $\Delta$ is one nilpotent Jordan block, even after scalar
extension to an algebraic closure of $\F_7$.
\end{theorem}
\begin{proof}
The polynomials $b_r$ have distinct degrees $r<7$ and nonzero leading
coefficients. Evaluation on the seven elements of $\F_7$ is injective for
polynomials of degree at most six, so these functions form a basis.
Pascal's polynomial identity gives $b_r(k+1)-b_r(k)=b_{r-1}(k)$, including
the closing edge $6\mapsto0$ because it is an identity in $\F_7[k]$.
This proves the Jordan chain and the rank assertions. Alternatively,
$\Delta^7=(T-I)^7=T^7-I=0$; its first power is nonzero, so a diagonal form
is impossible over any extension field.
\end{proof}

In this basis $\im\Delta=\langle b_0,\ldots,b_5\rangle$.
If $g=\sum_{r=0}^5c_rb_r$, a primitive is
$\sum_{r=0}^5c_rb_{r+1}$, plus a constant. This is a concrete alternative to
Fourier division. For instance $g=1=b_0$ has the primitive $b_1(k)=k$;
both are zero-sum as functions on $\F_7$, and adding any constant preserves
zero sum. Conversely $g=b_5$ has only primitives $b_6+c$.
Since $b_6$ has values $(0,0,0,0,0,0,1)$, no such primitive has zero sum.
Both examples satisfy $\sigma(g)=0$.

Over $\C$, a seven-cycle \emph{does} have a Fourier eigenbasis with seventh
roots of unity. The distinction here comes from the characteristic of the
coefficient field, not from the finite orbit alone.

\begin{center}
\begin{tabular}{@{}llll@{}}
\toprule
motion & moving orbit & characteristic-seven structure & normalization\\
\midrule
nonsplit $R$ & $8$ & Fourier over $\F_{49}$ & unique zero-sum primitive\\
split $S$ & $6$ & Fourier over $\F_7$ & unique zero-sum primitive\\
unipotent $U$ & $7$ & one nilpotent block & choose a base point\\
\bottomrule
\end{tabular}
\end{center}

\section{Covariance and the choice of one step}\label{sec:covariance}
On functions on $X=\PP^1(\F_7)$ let $T_gf(x)=f(gx)$,
$\Delta_g=T_g-I$, and $\mathcal U_hf(x)=f(h^{-1}x)$.
With this forward-shift convention, $T_aT_b=T_{ba}$.
\begin{proposition}[Transport of a difference]\label{prop:covariance}
For every $g,h\in\PGL_2(\F_7)$,
\[
 \mathcal U_h\Delta_g\mathcal U_h^{-1}=\Delta_{hgh^{-1}}.
\]
The stabilizer of $\Delta_R$ under this conjugation is $\langle R\rangle$,
of order eight. Its normalizer as a cyclic subgroup is dihedral of order
sixteen; the reflection $H(t)=-t$ sends $R$ to $R^{-1}$.
Thus $L_R=T_R+T_{R^{-1}}-2I$ is invariant under this normalizer, whereas
$\Delta_R$ is not.
\end{proposition}
\begin{proof}
Substitute $h^{-1}x$ into $T_g\mathcal U_h^{-1}f$. Equality of the difference
operators is equality of the corresponding permutations, since point
indicators distinguish them. Every permutation commuting with an eight-cycle
is a power of it, which proves the centralizer assertion. Over $\F_{49}$,
the nonsplit torus has exactly two common fixed points; any element
normalizing it must preserve this unordered pair. Fixing both points gives
the diagonal torus in that eigenbasis. Its $\F_7$-rational part is the
eight-element group in Proposition~\ref{prop:torus}: a rational linear map
preserving both eigenlines commutes with their generating matrix and is of
the form $aI+b\left(\begin{smallmatrix}0&1\\-1&0\end{smallmatrix}\right)$.
The map $H$ exchanges the pair and conjugates the torus by inversion.
This gives precisely two cosets and proves the normalizer assertion.
The last statements follow by exchanging $R$ and $R^{-1}$.
\end{proof}

Even retaining the torus but forgetting its generator is insufficient to
select $L_R$: $R^3$ is also a generator, and $L_{R^3}\ne L_R$.
For example their off-diagonal supports at a point are respectively its
neighbors at steps $\pm1$ and $\pm3$ on the eight-cycle. More generally
\[
 \Delta_{R^a}=(I+T_R+\cdots+T_R^{a-1})\Delta_R\quad(a\ge1),
\]
so a different step changes the operator explicitly.

\begin{proposition}[Full projective invariance is much more restrictive]
\label{prop:commutant}
Over any coefficient field $E$, every linear operator on $E^X$ commuting
with the full $\PGL_2(\F_7)$-action is of the form
\[
 f\longmapsto a f+b\,\sigma(f)\mathbf1\qquad(a,b\in E).
\]
In particular a directed one-step difference is not intrinsic to the bare
projective conic.
\end{proposition}
\begin{proof}
The group is two-transitive on $X$: it is transitive on ordered pairs of
distinct lines in $\F_7^2$ by changing a basis. A commuting operator's matrix
entries must therefore be constant on the diagonal and constant off it.
This is exactly the displayed form. An eight-cycle permutation matrix
minus $I$ does not have constant off-diagonal entries.
\end{proof}

\section{A real elliptic flow and a compactness obstruction}
The following assertions are over $\R$ and are independent of reduction
modulo seven. For homogeneous parameters $[s_1:s_0]\in\PP^1(\R)$ put
\[
 R_s=\left[\begin{pmatrix}s_0&s_1\\-s_1&s_0\end{pmatrix}\right].
\]
All these matrices are invertible. Multiplication gives
\begin{equation}\label{eq:real-law}
 [s_1:s_0]\star[r_1:r_0]
  =[s_1r_0+s_0r_1:s_0r_0-s_1r_1].
\end{equation}
For finite parameters this is $(s+r)/(1-sr)$, with its homogeneous meaning
when $sr=1$. For $s\ne0$ the projective map is the four-context history
\begin{equation}\label{eq:real-word}
 R_s=A_{-1/s}\circ M_{-(1+s^2)/s^2}\circ J\circ A_{-1/s}.
\end{equation}
The only ordinary input exclusion of this word is $t=1/s$.
Formula~\eqref{eq:real-word} also holds over any field of characteristic
different from two with $s\ne0$ and $1+s^2\ne0$.

\begin{proposition}[Global projective elliptic flow]\label{prop:real-flow}
The matrices
\[
 G_h=\begin{pmatrix}\cos h&\sin h\\-\sin h&\cos h\end{pmatrix}
\]
induce a smooth flow on $\PP^1(\R)$ of period $\pi$. In its finite chart,
\[
 \Phi_h(t)=\frac{t\cos h+\sin h}{\cos h-t\sin h},\qquad
 \frac{dt}{dh}=1+t^2.
\]
Its induced motion on the conic under $\nu$ rotates through angle $2h$.
The affine-coordinate blow-up is a chart crossing, not a failure of the
projective flow. It remains a pole for ordinary arithmetic in that chart.
\end{proposition}
\begin{proof}
Matrix multiplication gives $G_hG_k=G_{h+k}$. The first positive scalar
matrix occurs at $h=\pi$, proving the projective period. Differentiating at
$h=0$ gives $1+t^2$. In the other chart $u=1/t$, its vector field is
$-(1+u^2)\partial_u$, which is smooth at $u=0$. In the conic coordinate
$w(t)=(1+\ii t)/(1-\ii t)$ the multiplier is $e^{2\ii h}$ by
Proposition~\ref{prop:torus}, proving the double angle.
\end{proof}

These are elliptic motions within the projective group already generated by
Paper I's histories. They do not supply a two-dimensional regular AES.
\begin{proposition}[Compactness obstruction for the existing AES equation]
\label{prop:compact}
Let $(M,g)$ be a nonempty compact smooth Riemannian manifold without boundary,
and let $a:M\to\R$ be smooth. There are no real constants $\mu\ne0$, $\lambda$
for which
\[
 |\nabla a|_g^2=\mu^2+\lambda^2a^2
\]
holds everywhere.
\end{proposition}
\begin{proof}
At a maximum of $a$, compactness and absence of boundary give $da=0$.
The left side is zero there and the right side is at least $\mu^2>0$.
\end{proof}
An extension to a compact positive-curvature surface must therefore specify
what changes: domain, assignment object, equation, or regularity category.
The one-dimensional projective flow proves none of those choices by itself.

\section{Paper I interface, verification, and next problems}
This note is maintained in the AEG research notes. The geometric foundation
belongs to Paper 0 under the September 22, 2026
\href{../../governance/00b-paper-0-geometric-foundation-amendment.md}{geometric-foundation amendment}.
Paper I's \href{../../paper-1/sections/03-projective-affine.tex}{projective and affine section} contains the
arbitrary-field generation theorem and states the affine limitation of its
flow theory. Source provenance and migration details are recorded in the
companion resource directory's README.

Theorem~\ref{thm:history-witness} is a self-contained candidate example after
the projective-evaluation/domain warning in Paper I. It adds an explicit
nonsplit finite cyclic history, rather than strengthening the generation
theorem. Remark~\ref{rem:double} is the necessary coordinate/representation
warning. The finite difference theory belongs in these research notes until
an explicit analytic scope change is made. The real elliptic flow is a proved
research interface, not a claim that Paper I's affine flow has already been
extended as an AES theory. No AEG manuscript or governance status is changed
by this note.

The exact script in the companion \texttt{conic-dynamics} resource directory enumerates all $336$ projective matrices
over $\F_7$ independently of the circle parametrization. It checks the
four-context history and every ordinary input domain; the quotient/norm-one
intertwiner; the three orbit decompositions; difference ranks and Jordan
chains; the Fourier formulas on basis vectors, including Frobenius descent;
and covariance against every projective matrix on every point indicator.
It also checks normalizer and centralizer sizes and the corrected distinction
between nonsquare directions and isotropic directions. A successful finite
check does not replace the general-field proofs above.

Three next questions remain open as research tasks:
\begin{enumerate}
\item Extend the proved three-cycle comparison to general odd $q=p^r$.
The torus orders are $q-1$ and $q+1$, but an individual nontrivial translation
has order $p$, not $q$. Thus for $r>1$ its affine orbits are $p$-cycles, and
there is no single translation step visiting every affine point.
\item Compare histories producing the same projective torus element while
retaining their ordinary domains and marked context words. Determine what
additional process information a proposed invariant actually retains.
\item Formulate an elliptic AES category that states its assignment, domain,
critical set, and relation to the existing affine model. Prove any process or
curvature comparison separately from the projective one-dimensional results.
\end{enumerate}
Gauss sums, local-field trees, and group-algebra expansions are deferred until
one of these precise questions needs them. In particular, a finite projective
group is not by itself the acting group of a Bruhat--Tits tree; that step
requires a specified non-Archimedean local field and a reduction comparison.
\end{document}
